% GENERATED FILE. Edit canonical lessons, then run npm run generate:books.
% canonical-source-hash: fnv1a64:d3ce5937b41c7336
% canonical-lessons: JA-C160-souji, JA-C160-chiri, JA-C160-yogoreru, JA-C160-yogosu, JA-C160-kitanai

\chapter{Cleaning Up Dust and Dirt}
\label{ch:ja-cleaning-up-dust-and-dirt}
\hlchaptermodality{160}

\begin{chapteropening}
\textbf{By the end of this chapter:} \emph{I can say what gets dirty and how I clean it.}

Complete the last lesson of chapter 160: I can say what gets dirty and how I clean it.
\end{chapteropening}

\section[\texorpdfstring{\ja{そうじ} (sōji)}{そうじ (sōji)}]{\ja{そうじ} (sōji) — cleaning; sweeping}

\label{lesson:JA-C160-souji}



\begin{quote}
Before the new one: say the Japanese for to stack, then the Japanese for sorting out.
\end{quote}

\subsection*{You'll want to know: \ja{そうじ}}
\textbf{\ja{そうじ}} — \emph{sōji} — \textquotedblleft{}cleaning; sweeping\textquotedblright{}.

\begin{quote}
\textbf{\ja{そうじは} \ja{どようびです。}} — \emph{sōji wa doyōbi desu} — \textquotedblleft{}The cleaning is on Saturdays.\textquotedblright{}
\end{quote}

\begin{cousinweb}
\textbf{\ja{そうじ}} (\emph{sōji}) is Sino-Japanese (kango): two Chinese characters read with their on'yomi, readings borrowed from Chinese. They are \emph{sō}, \textquotedblleft{}sweep\textquotedblright{} (Middle Chinese \emph{sawX}) and \emph{ji}, \textquotedblleft{}remove\textquotedblright{} (\emph{drjo}) (KANJIDIC2; Middle Chinese in Baxter's transcription).
\end{cousinweb}

\begin{grammarlens}[title={What you've built}]
One more word for everyday life at A2.
\end{grammarlens}

\subsection*{Your turn}
\begin{itemize}
\raggedright
  \item \emph{Say it:} \emph{sōji}
  \item \emph{Say it:} \emph{sōji}, once more
  \item \emph{Say it:} \emph{sōji wa doyōbi desu} — The cleaning is on Saturdays.
  \item \emph{Recall:} say the Japanese for to stack, then the Japanese for sorting out, then say \emph{sōji} again
\end{itemize}

\begin{tcolorbox}[breakable,colback=teal!4,colframe=teal!35!black,title={Before you move on}]
What does \ja{そうじ} mean? (\textquotedblleft{}Cleaning; sweeping\textquotedblright{}.) Say it once more.
\end{tcolorbox}
\section[\texorpdfstring{\ja{ちり} (chiri)}{ちり (chiri)}]{\ja{ちり} (chiri) — dust specks; litter}

\label{lesson:JA-C160-chiri}



\begin{quote}
Before the new one: say the Japanese for sorting out, then the Japanese for cleaning.
\end{quote}

\subsection*{You'll want to know: \ja{ちり}}
\textbf{\ja{ちり}} — \emph{chiri} — \textquotedblleft{}dust specks; litter\textquotedblright{}.

\begin{quote}
\textbf{\ja{ちりとりで} \ja{ちりを} \ja{ひろいます。}} — \emph{chiritori de chiri o hiroimasu} — \textquotedblleft{}I pick up the dust with a dustpan.\textquotedblright{}
\end{quote}

\begin{cousinweb}
\textbf{\ja{ちり}} (\emph{chiri}) is a native Japanese word (wago), inherited rather than borrowed from Chinese, and its older origin is not settled; the kanji JMdict gives for it was attached later (JMdict).
\end{cousinweb}

\begin{grammarlens}[title={What you've built}]
One more word for everyday life at A2.
\end{grammarlens}

\subsection*{Your turn}
\begin{itemize}
\raggedright
  \item \emph{Say it:} \emph{chiri}
  \item \emph{Say it:} \emph{chiri}, once more
  \item \emph{Say it:} \emph{chiritori de chiri o hiroimasu} — I pick up the dust with a dustpan.
  \item \emph{Recall:} say the Japanese for sorting out, then the Japanese for cleaning, then say \emph{chiri} again
\end{itemize}

\begin{tcolorbox}[breakable,colback=teal!4,colframe=teal!35!black,title={Before you move on}]
What does \ja{ちり} mean? (\textquotedblleft{}Dust specks; litter\textquotedblright{}.) Say it once more.
\end{tcolorbox}
\section[\texorpdfstring{\ja{よごれる} (yogoreru)}{よごれる (yogoreru)}]{\ja{よごれる} (yogoreru) — to get dirty}

\label{lesson:JA-C160-yogoreru}



\begin{quote}
Before the new one: say the Japanese for cleaning, then the Japanese for dust specks.
\end{quote}

\subsection*{You'll want to know: \ja{よごれる}}
\textbf{\ja{よごれる}} — \emph{yogoreru} — \textquotedblleft{}to get dirty\textquotedblright{}.

\begin{quote}
\textbf{\ja{こどもの} \ja{ふくは} \ja{よく} \ja{よごれます。}} — \emph{kodomo no fuku wa yoku yogoremasu} — \textquotedblleft{}Children's clothes often get dirty.\textquotedblright{}
\end{quote}

Polite \textbf{\ja{よごれます}}, \emph{yogoremasu}; an ichidan verb.

\begin{cousinweb}
\textbf{\ja{よごれる}} (\emph{yogoreru}) is a native Japanese word (wago), inherited rather than borrowed from Chinese; its kanji was attached later as a kun'yomi, and its older origin is not settled (JMdict, KANJIDIC2).
\end{cousinweb}

\begin{grammarlens}[title={What you've built}]
One more word for everyday life at A2.
\end{grammarlens}

\subsection*{Your turn}
\begin{itemize}
\raggedright
  \item \emph{Say it:} \emph{yogoreru}
  \item \emph{Say it:} \emph{yogoreru}, once more
  \item \emph{Say it:} \emph{kodomo no fuku wa yoku yogoremasu} — Children's clothes often get dirty.
  \item \emph{Recall:} say the Japanese for cleaning, then the Japanese for dust specks, then say \emph{yogoreru} again
\end{itemize}

\begin{tcolorbox}[breakable,colback=teal!4,colframe=teal!35!black,title={Before you move on}]
What does \ja{よごれる} mean? (\textquotedblleft{}To get dirty\textquotedblright{}.) Say it once more.
\end{tcolorbox}
\section[\texorpdfstring{\ja{よごす} (yogosu)}{よごす (yogosu)}]{\ja{よごす} (yogosu) — to make dirty; to soil}

\label{lesson:JA-C160-yogosu}



\begin{quote}
Before the new one: say the Japanese for dust specks, then the Japanese for to get dirty.
\end{quote}

\subsection*{You'll want to know: \ja{よごす}}
\textbf{\ja{よごす}} — \emph{yogosu} — \textquotedblleft{}to make dirty; to soil\textquotedblright{}.

\begin{quote}
\textbf{\ja{いぬが} \ja{ゆかを} \ja{よごします。}} — \emph{inu ga yuka o yogoshimasu} — \textquotedblleft{}The dog makes the floor dirty.\textquotedblright{}
\end{quote}

Polite \textbf{\ja{よごします}}, \emph{yogoshimasu}; a godan verb. Its intransitive partner is \textbf{\ja{よごれる}}, \emph{yogoreru}.

\begin{cousinweb}
\textbf{\ja{よごす}} (\emph{yogosu}) is a native Japanese word (wago), inherited rather than borrowed from Chinese, and built on the same root as \emph{yogoreru}, \textquotedblleft{}to get dirty\textquotedblright{}; its kanji was attached later as a kun'yomi (JMdict, KANJIDIC2).
\end{cousinweb}

\begin{grammarlens}[title={What you've built}]
One more word for everyday life at A2.
\end{grammarlens}

\subsection*{Your turn}
\begin{itemize}
\raggedright
  \item \emph{Say it:} \emph{yogosu}
  \item \emph{Say it:} \emph{yogosu}, once more
  \item \emph{Say it:} \emph{inu ga yuka o yogoshimasu} — The dog makes the floor dirty.
  \item \emph{Recall:} say the Japanese for dust specks, then the Japanese for to get dirty, then say \emph{yogosu} again
\end{itemize}

\begin{tcolorbox}[breakable,colback=teal!4,colframe=teal!35!black,title={Before you move on}]
What does \ja{よごす} mean? (\textquotedblleft{}To make dirty; to soil\textquotedblright{}.) Say it once more.
\end{tcolorbox}
\section[\texorpdfstring{\ja{きたない} (kitanai)}{きたない (kitanai)}]{\ja{きたない} (kitanai) — dirty; messy}

\label{lesson:JA-C160-kitanai}



\begin{quote}
Before the new one: say the Japanese for to get dirty, then the Japanese for to make dirty.
\end{quote}

\subsection*{You'll want to know: \ja{きたない}}
\textbf{\ja{きたない}} — \emph{kitanai} — \textquotedblleft{}dirty; messy\textquotedblright{}.

\begin{quote}
\textbf{\ja{この} \ja{へやは} \ja{きたないです。}} — \emph{kono heya wa kitanai desu} — \textquotedblleft{}This room is dirty.\textquotedblright{}
\end{quote}

\begin{cousinweb}
\textbf{\ja{きたない}} (\emph{kitanai}) is a native Japanese word (wago), inherited rather than borrowed from Chinese; its kanji was attached later as a kun'yomi, and its older origin is not settled (JMdict, KANJIDIC2).
\end{cousinweb}

\begin{grammarlens}[title={What you've built}]
One more word for everyday life at A2.
\end{grammarlens}

\subsection*{Your turn}
\begin{itemize}
\raggedright
  \item \emph{Say it:} \emph{kitanai}
  \item \emph{Say it:} \emph{kitanai}, once more
  \item \emph{Say it:} \emph{kono heya wa kitanai desu} — This room is dirty.
  \item \emph{Recall:} say the Japanese for to get dirty, then the Japanese for to make dirty, then say \emph{kitanai} again
\end{itemize}

\begin{tcolorbox}[breakable,colback=teal!4,colframe=teal!35!black,title={Before you move on}]
What does \ja{きたない} mean? (\textquotedblleft{}Dirty; messy\textquotedblright{}.) Say it once more.
\end{tcolorbox}
